\documentclass[french]{article}
\usepackage[utf8]{inputenc}
\usepackage[T1]{fontenc}
\usepackage{lmodern}
\usepackage{geometry}
\usepackage{graphicx}
\usepackage{babel}
\usepackage{amsmath}
\usepackage{amsthm}
\usepackage{amssymb}
\usepackage{hyperref}
\usepackage{stmaryrd}
\usepackage{enumitem}
\usepackage{tcolorbox}
\usepackage{dsfont}
\usepackage{multirow}
\usepackage{etoc}
\usepackage{titlesec}
\usepackage{esint}
\usepackage{cancel}
\usepackage{mathdots}

\setlist[itemize]{label=$\bullet$}


\renewcommand{\P}{\mathbb{P}}
\newcommand{\N}{\mathbb{N}}
\newcommand{\Z}{\mathbb{Z}}
\newcommand{\Q}{\mathbb{Q}}
\newcommand{\R}{\mathbb{R}}
\newcommand{\C}{\mathbb{C}}
\newcommand{\K}{\mathbb{K}}
\renewcommand{\L}{\mathbb{L}}
\newcommand{\U}{\mathbb{U}}
\newcommand{\st}{^*}
\newcommand{\tq}{\middle|}
\newcommand{\inv}{^{-1}}
\newcommand{\E}{\mathbb{E}}
\newcommand{\F}{\mathbb{F}}
\newcommand{\V}{\mathbb{V}}
\renewcommand{\ker}{\text{Ker}}
\newcommand{\im}{\text{Im}}
\newcommand{\card}{\text{Card}}
\newcommand{\Tr}{\text{Tr}}

\title{TD Polynômes \\ Exercice 12}
\date{}
\begin{document}
\maketitle

\section{Énoncé}
Soit $P\in\R[X]$ unitaire de degré $n\geqslant1$. Montrer que $P$ est scindé sur $\R$ si, et seulement si, $$\forall z\in\C,\ \lvert P(z)\rvert\geqslant\lvert\im(z)\rvert^n$$
\section{Solution}
On procède par double implication :
\begin{itemize}
	\item Sens direct : si $P$ est scindé sur $\R$, on écrit $$P=\prod_{k=1}^{n}(X-\lambda_k)$$ Alors pour tout $z\in\C$ : $$\lvert P(z)\rvert=\prod_{k=1}^{n}\lvert z-\lambda_k\rvert \geqslant \prod_{k=1}^{n}\lvert\im(z-\lambda_k)\rvert=\prod_{k=1}^{n}\lvert\im(z)\rvert=\lvert\im(z)\rvert ^n$$
	\item Sens réciproque : notons $\lambda_k$ les racines \textit{a priori} complexes de $P$ après l'avoir scindé sur $\C$. On a alors : $$\lvert\im(\lambda_k)\rvert^n\leqslant\lvert P(\lambda_k)\rvert=0$$ Puisque $n\geqslant1$, on en déduit que $\im(\lambda_k)=0$ et donc que $\lambda_k\in\R$, si bien que $P$ est scindé sur $\R$.
\end{itemize}
\end{document}
